\section{Implementation and data details}\label{app:impl}

\paragraph{Encoder.}
The frozen checkpoint is V-JEPA\,2.1 ViT-L distilled from ViT-G at
$384$ pixels. Its input is bilinearly resized to $256\times256$ and
normalised using ImageNet statistics. Tubelet size is two, patch size is
$16$, and the $24$-block encoder is read at zero-indexed blocks
$5,11,17,23$ after the corresponding hierarchical normalisation.
This produces four $2048\times1024$ fields.
The frozen encoder contributes approximately $304.7$\,M parameters in
addition to the $22.4$\,M trainable pixel branch.
Parameter count of the pixel branch alone is not an inference-cost comparison.

\paragraph{Pixel branch.}
Attention heads are $(2,4,4,8,8)$ and the feed-forward expansion is $2.66$.
The decoder mirrors the four non-bottleneck encoder scales, with skip
concatenation and a $1\times1$ reduction.
Pixel-unshuffle and pixel-shuffle operations change resolution.
The bottleneck has a residual temporal convolution followed by cross-frame
attention; other restoration blocks operate frame-wise.
Each stage applies restoration blocks, token cross-attention where enabled,
and local illumination cue injection.
Reflectance and illumination heads are convolution--SiLU--convolution--sigmoid
stacks with three and one output channels, respectively. Their outputs
contribute only to the training loss.
The residual head is a zero-initialised $3\times3$ convolution.

\paragraph{Initialisation and optimisation.}\label{app:stages}
The converged CVPR comparisons use from-scratch pixel branches, not the
warm-started stages in the earlier project snapshot.
The frozen representation is the only pretrained part of the default model.
The zero-token arm keeps the same cross-attention modules.
The arms use the same nominal seed set, but constructing the pretrained
encoder consumes random draws before the pixel branch is initialised,
whereas the zero-token arm skips that construction. Their initial pixel
weights and training random draws are therefore not shared.
The comparison controls architecture and optimisation recipe; pairing refers
to test sequences, not to identical training randomness.
All arms use the optimiser and crop recipe in Sec.~\ref{sec:method:train}
and score the final checkpoint without validation-based selection.
The converged run collection includes both Ascend 910B3 and NVIDIA A100
runs; experiment identifiers, seeds and hardware are listed in the released
manifest. The third in-domain seed was completed on different hardware, so
the seed collection does not separately estimate hardware effects.

\paragraph{Split provenance.}\label{app:membership}
SDSD-indoor contributes $46/6/14$ train/validation/test sequences and
SDSD-outdoor $90/10/25$.
Validation groups repeated passes of a scene together.
BVI-RLV contains both frame rates of each scene; its two clips per scene
stay in the same data split.
LOL-v1 has $436/49/15$ training/validation/test images,
LOL-v2-real $620/69/100$, and LOL-v2-synthetic $810/90/100$.
The frozen split files, rather than a split reconstructed at evaluation time,
define membership.
Retinexformer is retrained on these SDSD training sequences.
Results from released checkpoints trained on other splits are excluded from
the current main comparison.

\paragraph{Evaluation details.}
Inputs and targets are paired by sorted frame index, which handles SDSD
sequences with filename offsets.
The central video clip and spatial crop are fixed before inference.
Short clips and images are repeated to $16$ frames; images produce a single
scored frame.
Image sizes are padded to the branch's spatial divisor and cropped back
after inference.
Raw predictions are clipped to $[0,1]$ and rounded to uint8 for PSNR.
GT-mean rescaling and previous colour-calibration procedures are not used in
the main claims.
The registered bootstrap unit is the evaluated sequence or image.
Paired SDSD passes and BVI-RLV frame-rate clips can share a scene.
A post-hoc sensitivity analysis therefore resamples whole scene groups,
retaining sequence weights through a total-sum/total-count estimator.
The in-domain PSNR-Y and LPIPS benefits and the two video-domain transfer
LPIPS benefits retain their direction and zero-excluding intervals.
The SDSD transfer LPIPS bound is close to zero. This robustness analysis
does not replace more training seeds or eliminate other dependence.

\section{Intervention details}\label{app:interventions}

\paragraph{Principal components.}
Per-clip projection centres all $2048$ tokens at each tap, obtains the right
singular vectors of the centred matrix, and reconstructs it after retaining
$k$ directions. The full clip mean is added back.
It therefore retains substantially more than $k$ scalar values: the mean,
directions, and token-dependent coordinates all remain.
The fixed basis uses the first $60$ training-split clips in the saved order,
accumulating means and second moments separately at each tap; eigendecomposition
gives a rank-$64$ basis, truncated to the requested rank.
No target or test label is used in either projection.
The per-clip projection nevertheless depends on the test input, as intended,
whereas the fixed basis and mean do not.

\paragraph{Exploratory mean--basis control (validation only).}\label{app:meanbasis}
The historical test intervention changes both factors.
To separate them, we use the existing full-token-trained E17 seeds zero and
one on every frozen SDSD validation sequence, without retraining or checkpoint
selection. The diagnostic was chosen after the earlier results were inspected;
its six conditions and validation-only protocol were recorded in Amendment~21
before its outcomes. It is exploratory rather than a new confirmatory test.
All principal-component projections here have rank four.
The adaptive directions are computed from clip-centred tokens independently
of which mean is used in reconstruction.

\begin{table}[t]
\centering\small
\setlength{\tabcolsep}{4pt}
\begin{tabular}{lcc}
\toprule
Validation condition & PSNR-Y $\uparrow$ & LPIPS $\downarrow$\\
\midrule
Identity & \R{meanval@identity@psnr} & \R{meanval@identity@lpips}\\
Clip mean only & \R{meanval@meanonly@psnr} & \R{meanval@meanonly@lpips}\\
Clip mean, clip basis & \R{meanval@clipclip@psnr} & \R{meanval@clipclip@lpips}\\
Clip mean, fixed basis & \R{meanval@clipfixed@psnr} & \R{meanval@clipfixed@lpips}\\
Train mean, clip basis & \R{meanval@trainclip@psnr} & \R{meanval@trainclip@lpips}\\
Train mean, fixed basis & \R{meanval@trainfixed@psnr} & \R{meanval@trainfixed@lpips}\\
\bottomrule
\end{tabular}
\caption{Exploratory mean--basis diagnostic on
$\R{meanval@n}$ validation sequences from $\R{meanval@scenes}$ scenes,
averaged over $\R{meanval@seeds}$ existing checkpoints. The four crossed
conditions use rank-four directions; clip-mean-only broadcasts the mean to
every token. These are validation measurements, separate from
Table~\ref{tab:sufficiency}'s test results.}
\label{tab:meanbasis}
\end{table}

We pair differences by sequence after seed averaging and report unadjusted
$95\%$ scene-bootstrap intervals, with $10{,}000$ resamples from eight scenes;
the archive also reports sequence-bootstrap intervals and every seed effect.
At fixed clip mean, replacing the training basis with adaptive directions
changes PSNR-Y by $\R{meanval@clipvsfixed@psnr}$\,dB
($\R{meanval@clipvsfixed@psnr@sceneci}$).
At fixed adaptive basis, replacing the training mean with the clip mean
changes it by $\R{meanval@meanvsfixedmean_clipbasis@psnr}$\,dB
($\R{meanval@meanvsfixedmean_clipbasis@psnr@sceneci}$);
with the fixed basis, the corresponding change is
$\R{meanval@meanvsfixedmean_fixedbasis@psnr}$\,dB
($\R{meanval@meanvsfixedmean_fixedbasis@psnr@sceneci}$).
Every seed agrees in sign on these contrasts.
The basis contrast at the training mean remains unresolved on PSNR,
so the adaptive-basis effect is conditional on the retained mean.

Rank-four adaptive projection also improves on clip-mean-only by
$\R{meanval@rank4vsmeanonly@psnr}$\,dB
($\R{meanval@rank4vsmeanonly@psnr@sceneci}$), with both seeds positive.
Against identity, however, the change is
$\R{meanval@rank4vsidentity@psnr}$\,dB
($\R{meanval@rank4vsidentity@psnr@sceneci}$), and LPIPS effects have
opposite signs across seeds. This does not establish equivalence or a
validation improvement over the full field.
The matched contrasts support contributions from both the input-dependent
mean and basis under these interventions. They do not show that either can be
obtained cheaply, that a trained model needs only four dimensions, or that the
same representation suffices during training.

\paragraph{Interpretation of averaging.}
Spatial mean retains a separate vector for each tubelet; temporal mean retains
a separate vector for each spatial position.
Neither transformation removes the frozen encoder's computation.
An interval that crosses zero bounds the detected effect but does not establish
statistical equivalence without a predeclared equivalence margin.

\paragraph{Temporal measurements.}
The latent diagnostic reports mean cosine agreement of consecutive-tubelet
differences and their mean absolute magnitude ratio.
The relative and cosine losses are trained on two seeds each.
The cosine form's log-ratio penalty penalises both expansion and contraction.
For each sequence and arm, tLP and target tLP are averaged over the arm's
seeds, then their absolute difference is taken. Differences between these
absolute gaps are paired over sequences. Each individual seed's absolute-gap
effect is also checked for the same sign.
The order of averaging and taking the absolute value matches the released
aggregation tool. This gap uses identical frames and crops and is not a
motion-compensated error.

\section{Corrected controls and outstanding validation}\label{app:pending}

The archive contains evaluations beyond those used in the current claims.
The random-encoder control has now been rescored with CPU initialisation seed
$1234$, matching the training implementation. Its recorded initial-parameter
checksum agrees up to the float-versus-double accumulation difference logged
by the NPU runtime; this is an approximate fingerprint, not a full model hash.
The corrected JSON records reconstruction settings separately from the
historical result. Section~\ref{sec:exp:control} uses the corrected single-seed
comparison and does not retain the earlier large fidelity-attribution claim.

Corrected hybrid cross-frame keys are reported in
Sec.~\ref{sec:exp:ablation}. Restoring all-depth keys for the available
seed-two checkpoint gives $\R{audit@alllayers@psnr@a}$\,dB, replacing the
archived, incorrectly configured $21.568$\,dB result. Against the corresponding
full-token default seed, its differences are
$\R{audit@alllayers@psnr}$\,dB ($\R{audit@alllayers@psnr@ci}$) and
$\R{audit@alllayers@lpips}$ LPIPS ($\R{audit@alllayers@lpips@ci}$).
Both intervals include zero. The variant was trained on Ascend and the
reference on A100, with only one available variant seed; this conditional
comparison cannot establish a general architectural effect. Its temporal
diagnostics also show tradeoffs, detailed in the audit record.
Deepest-only taps and learned gates still require the actual trained
configuration and weights before their archived differences can support
architectural conclusions.
The fixed-rank-$64$ training control has been reevaluated with its saved
projection. Its current comparison uses the same two baseline seed numbers;
the old three-baseline-seed comparison is preserved in the audit record.
Neither comparison resolves a training-time penalty after correction.
It still does not establish a computational compression result.

H20/E28, training from scratch with per-clip centred rank-four tokens, remains
unfinished: no verified final evaluation is available.
It is neither a positive nor a negative result.
The proposed small-network distillation experiment was cancelled by its
registered screening rule and was not run; its cancellation supplies no
evidence that another encoder or distillation procedure cannot work.
Some baseline temporal diagnostics were not included in the original archive.
Claims requiring those missing measurements are withheld until they are
reproduced under the same protocol. The qualitative panels have been
regenerated with checkpoint, source and output hashes.

\section{Scope of the conclusions}\label{app:scope}

The main conclusions concern one frozen encoder, one restoration family,
and the specified data splits.
Few seeds limit inference about training variability, and sequence bootstraps
do not replace additional independent runs.
SDSD includes low-motion material, while the LOL test contains still images.
A video encoder's temporal sensitivity and a readout's weak order sensitivity
therefore do not prove that temporal representations are irrelevant to other
video-restoration tasks.
The tested low-rank replacements describe the behaviour of existing
checkpoints. The exploratory validation control separates mean and basis
effects under those replacements; future work must separately test efficient
encoders and training with the restricted representation.
